% !TeX root=../main.tex
% در این فایل، عنوان پایان‌نامه، مشخصات خود، متن تقدیمی‌، ستایش، سپاس‌گزاری و چکیده پایان‌نامه را به فارسی، وارد کنید.
% توجه داشته باشید که جدول حاوی مشخصات پروژه/پایان‌نامه/رساله و همچنین، مشخصات داخل آن، به طور خودکار، درج می‌شود.
%%%%%%%%%%%%%%%%%%%%%%%%%%%%%%%%%%%%
% دانشگاه خود را وارد کنید
\university{دانشگاه تهران}
% پردیس دانشگاهی خود را اگر نیاز است وارد کنید (مثال: فنی، علوم پایه، علوم انسانی و ...)
\college{پردیس دانشکده‌های فنی}
% دانشکده، آموزشکده و یا پژوهشکده  خود را وارد کنید
\faculty{دانشکده مهندسی مکانیک}
% گروه آموزشی خود را وارد کنید (در صورت نیاز)
\department{گروه طراحی کاربردی}
% رشته تحصیلی خود را وارد کنید
\subject{مهندسی مکانیک}
% گرایش خود را وارد کنید
\field{طراحی کاربردی}
% عنوان پایان‌نامه را وارد کنید
\title{ناوبری خودران اسکوتر چهارچرخ محرک در محیط‌های پویا بر پایه یادگیری تقویتی}
% نام استاد(ان) راهنما را وارد کنید
\firstsupervisor{دکتر مسعود شریعت پناهی}
\firstsupervisorrank{استاد}
\secondsupervisor{دکتر محمد رضا حائری یزدی}
\secondsupervisorrank{استاد}
% نام استاد(دان) مشاور را وارد کنید. چنانچه استاد مشاور ندارید، دستورات پایین را غیرفعال کنید.
\firstadvisor{دکتر هادی قطان کاشانی}
\firstadvisorrank{استادیار}
%\secondadvisor{دکتر مشاور دوم}
% نام داوران داخلی و خارجی خود را وارد نمایید.
\internaljudge{دکتر داور داخلی}
\internaljudgerank{دانشیار}
\externaljudge{دکتر داور خارجی}
\externaljudgerank{دانشیار}
\externaljudgeuniversity{دانشگاه داور خارجی}
% نام نماینده کمیته تحصیلات تکمیلی در دانشکده \ گروه
\graduatedeputy{دکتر نماینده}
\graduatedeputyrank{دانشیار}
% نام دانشجو را وارد کنید
\name{محمد مهدی}
% نام خانوادگی دانشجو را وارد کنید
\surname{عابدی}
% شماره دانشجویی دانشجو را وارد کنید
\studentID{۸۱۰۶۰۱۱۵۵}
% تاریخ پایان‌نامه را وارد کنید
\thesisdate{تیر ۱۴۰۵}
% به صورت پیش‌فرض برای پایان‌نامه‌های کارشناسی تا دکترا به ترتیب از عبارات «پروژه»، «پایان‌نامه» و «رساله» استفاده می‌شود؛ اگر  نمی‌پسندید هر عنوانی را که مایلید در دستور زیر قرار داده و آنرا از حالت توضیح خارج کنید.
%\projectLabel{پایان‌نامه}

% به صورت پیش‌فرض برای عناوین مقاطع تحصیلی کارشناسی تا دکترا به ترتیب از عبارت «کارشناسی»، «کارشناسی ارشد» و «دکتری» استفاده می‌شود؛ اگر نمی‌پسندید هر عنوانی را که مایلید در دستور زیر قرار داده و آنرا از حالت توضیح خارج کنید.
%\degree{}
%%%%%%%%%%%%%%%%%%%%%%%%%%%%%%%%%%%%%%%%%%%%%%%%%%%%
%% پایان‌نامه خود را تقدیم کنید! %%
\dedication
{
{\Large تقدیم به:}\\
\begin{flushleft}{
	\huge
	همسرم\\
	\vspace{7mm}
	و\\
	\vspace{7mm}
	پدر و مادرم
}
\end{flushleft}
}
%% متن قدردانی %%
%% ترجیحا با توجه به ذوق و سلیقه خود متن قدردانی را تغییر دهید.
\acknowledgement{
سپاس خداوندگار حکیم را که با لطف بی‌کران خود، آدمی را به زیور عقل آراست.

در آغاز وظیفه‌  خود  می‌دانم از زحمات بی‌دریغ اساتید  راهنمای خود،  جناب آقای دکتر مسعود شریعت پناهی و دکتر محمدرضا حائری یزدی، صمیمانه تشکر و  قدردانی کنم که در طول انجام این پایان‌نامه با نهایت صبوری همواره راهنما و مشوق من بودند و قطعاً بدون راهنمایی‌های ارزنده‌ ایشان، این مجموعه به انجام نمی‌رسید.

از جناب آقای دکتر کاشانی که  زحمت مشاوره‌، بازبینی و تصحیح این پایان‌نامه را تقبل فرمودند کمال امتنان را دارم.

%از همکاری و مساعدت‌های دکتر ... مسئول تحصیلات تکمیلی و سایر کارکنان دانشکده بویژه سرکار خانم ... کمال تشکر را دارم.

با سپاس بی‌دریغ خدمت دوستان گران‌مایه‌ام، خانم‌ها ... و آقایان ... در آزمایشگاه ...، که با همفکری مرا صمیمانه و مشفقانه یاری داده‌اند.

و در پایان، بوسه می‌زنم بر دستان خداوندگاران مهر و مهربانی، پدر و مادر عزیزم و بعد از خدا، ستایش می‌کنم وجود مقدس‌شان را و تشکر می‌کنم از خانواده عزیزم به پاس عاطفه سرشار و گرمای امیدبخش وجودشان، که بهترین پشتیبان من بودند.
}
%%%%%%%%%%%%%%%%%%%%%%%%%%%%%%%%%%%%
%چکیده پایان‌نامه را وارد کنید
\fa-abstract{
ناوبری خودکار وسایل نقلیه زمینی در محیط‌های پویا، نیازمند یکپارچه‌سازی هم‌زمان ادراک محیط، تصمیم‌گیری حرکتی، اجتناب از موانع، تضمین ایمنی و اجرای دقیق فرمان‌ها است. این مسئله در اسکوترهای چهارچرخ محرک، به دلیل جرم و اینرسی نسبتاً زیاد، محدودیت توقف، ناهمسانی پاسخ موتورها، نویز حسگرها و حضور موانع متحرک، پیچیدگی بیشتری دارد. هدف این پژوهش، طراحی، پیاده‌سازی و ارزیابی یک چارچوب یادگیری‌محور برای هدایت خودکار اسکوتر چهارچرخ محرک به سمت هدف، در محیط‌های دارای موانع ایستا و متحرک است.

در چارچوب پیشنهادی، یک سیاست مبتنی بر الگوریتم \lr{LSTM-SAC} برای تولید پیوسته فرمان‌های سرعت خطی و زاویه‌ای توسعه داده شد. به‌منظور کاهش فاصله میان آموزش شبیه‌سازی و اجرای عملی، از ساختار معلم--دانشجو استفاده شد؛ به‌گونه‌ای که عامل معلم با اطلاعات دقیق محیط و عامل دانشجو با مشاهدات نزدیک‌تر به داده‌های حسگری آموزش دید. فضای مشاهده شامل اطلاعات نسبی هدف، وضعیت حرکتی اسکوتر، داده‌های حسگرهای فراصوت، فرمان قبلی، سطح خطر و ویژگی‌های زمانی چند مانع بود. همچنین، آموزش عامل‌ها به‌صورت مرحله‌ای و با افزایش تدریجی پیچیدگی محیط انجام شد. برای کاهش احتمال برخورد نیز یک سپر ایمنی بر پایه فاصله، حرکت نسبی موانع و زمان تخمینی تا برخورد در مسیر اعمال فرمان قرار گرفت.

در ارزیابی شبیه‌سازی، عامل دانشجو با استفاده از داده‌های حسگری و سپر ایمنی فعال، در پیچیده‌ترین سناریوی بررسی‌شده به نرخ موفقیت (96\%) و نرخ برخورد صفر دست یافت. همچنین، در ارزیابی عملی سامانه طی ۲۰ آزمون مستقل، نرخ موفقیت (80\%) و نرخ برخورد (5\%) به دست آمد. این نتایج ضمن تأیید امکان انتقال سیاست آموزش‌دیده به سامانه واقعی، نشان دادند که سپر ایمنی نقش مؤثری در کاهش برخوردها و افزایش پایداری عملکرد دارد.

نتایج نشان می‌دهند که ترکیب یادگیری تقویتی عمیق، آموزش معلم--دانشجو، ادراک چندحسگری و سپر ایمنی، امکان اجرای عملی ناوبری هدف‌محور اسکوتر را فراهم می‌کند. بااین‌حال، وابستگی سیاست به سپر ایمنی، خطاهای ادراک در حین حرکت، محدودیت کنترل موتورها و شکاف باقی‌مانده میان شبیه‌سازی و واقعیت، مهم‌ترین موانع دستیابی به عملکرد کاملاً پایدار و مستقل هستند.

}
% کلمات کلیدی پایان‌نامه را وارد کنید
\keywords{ناوبری خودکار، اسکوتر چهارچرخ محرک، یادگیری تقویتی عمیق}
% انتهای وارد کردن فیلد‌ها
%%%%%%%%%%%%%%%%%%%%%%%%%%%%%%%%%%%%%%%%%%%%%%%%%%%%%%
